\chapter{Grundlagen}

\section{Internal Audit}

Das Internal Audit stellt eine unabhängige und objektive Prüfungs- und Beratungsfunktion innerhalb einer Organisation dar. Die zentrale Aufgabe besteht darin, die Korrektheit und Wirksamkeit von Governance-, Risikomanagement- und internen Kontrollprozessen zu bewerten sowie durch unabhängige Prüfungs- und Beratungsleistungen zur Verbesserung dieser Prozesse beizutragen. Im Gegensatz zum External Audit ist das Internal Audit organisatorisch in das Unternehmen eingebunden, wobei seine unabhängige Position gegeben sein muss (vgl. \cite[S.~12~ff., S. 46]{IIA2024}).
Die Existenz von Internal-Audit-Funktionen begründet sich insbesondere durch die Notwendigkeit, die Transparenz organisatorischer Prozesse zu erhöhen, Rechenschaftspflichten gegenüber Stakeholdern zu unterstützen sowie das Vertrauen in Governance-, Risiko- und Kontrollprozesse zu stärken (vgl. \cite[S.~1]{IIA_BoardRole2023}).

Moderne Audit-Ansätze gehen dabei über die reine Überprüfung der Regelkonformität hinaus. Das Internal Audit soll vielmehr organisatorisches Lernen fördern, das Risikobewusstsein der beteiligten Akteure stärken und Verbesserungsmöglichkeiten in Geschäftsprozessen aufzeigen. Dadurch leistet die Funktion einen Beitrag zur kontinuierlichen Weiterentwicklung der Organisation und schafft einen Mehrwert, der über reine Compliance-Aspekte hinausgeht (vgl. \cite[S.~262 f.]{cantu2025making}). Für die Durchführung der Prüfungen werden unter anderem Prozess- und Ablaufdokumente gesammelt, interpretiert und analysiert sowie deren Übereinstimmung mit gegebenen Anforderungen bewertet (vgl. \cite[S.~556 f.]{ROUSSY2013550}). Gerade dieser dokumentenbasierte Teil der Prüfung bietet ein erhebliches Potenzial für den Einsatz von \ac{KI} (vgl. \cite[S.~1]{schumann2021natural}).

\section{Stichprobenanalyse im Internal Audit}

Für ein Prüfungsurteil muss das Internal Audit alle für das Prüfungsziel relevanten Informationen erheben (vgl. \cite[S.~104]{IIA2024}). Diese Informationen werden als Prüfungsnachweise bezeichnet (vgl. \cite[S.~5]{isa500}). Da Personal, Zeit und Budget begrenzt sind, kann das Internal Audit nicht alle Elemente prüfen und arbeitet deshalb mit Stichproben (vgl. \cite[S.~72]{vandernest2015sampling}). Eine Stichprobe liegt vor, wenn Prüfungshandlungen auf weniger als 100~\% der Elemente einer Grundgesamtheit angewendet werden, jedes Element ausgewählt werden kann und Schlüsse auf die gesamte Grundgesamtheit gezogen werden sollen (vgl. \cite[S.~12]{isa530}). Die Grundgesamtheit ist eine abgegrenzte Menge gleichartiger Elemente, etwa alle Rechnungen eines Kreditorenkontos (vgl. \cite[S.~12]{baumeister2019sample}). Zunächst wird eine Pflichtprobe aller Elemente bestimmt, die vorab festgelegte Kriterien erfüllen, etwa alle Rechnungen über einem bestimmten Betrag. Aus den übrigen Elementen wird anschließend eine Zufallsstichprobe gezogen (vgl. \cite{freudenberg2026auditsample}). Da nur ein Teil der Grundgesamtheit geprüft wird, entsteht ein Stichprobenrisiko. Dieses beschreibt das Risiko, dass auf Basis der Stichprobe ein anderes Ergebnis erzielt wird als bei einer Prüfung aller Elemente (vgl. \cite[S.~21]{isa530}). Eine Audit Exception liegt vor, wenn bei einem geprüften Element eine Abweichung von den festgelegten Anforderungen oder eine sonstige prüfungsrelevante Auffälligkeit festgestellt wird (vgl. \cite[S.~11]{IIA2024}).

\section{Large Language Models}
\label{sec:llm}

\acp{LLM} sind ein Teilbereich der generativen \ac{KI}. Unter generativer \ac{KI} werden Verfahren verstanden, die auf Grundlage von Trainingsdaten neue Inhalte wie Texte, Bilder oder Audio erzeugen können (vgl. \cite[S.~111]{Feuerriegel2024GenerativeAI}). Die Modelle werden mit einer großen Anzahl an Parametern auf umfangreichen Textdaten trainiert. Ihre Grundlage bildet die Sprachmodellierung, deren Ziel darin besteht, die Wahrscheinlichkeitsverteilung über Folgen von Wörtern einer Sprache zu erlernen (vgl. \cite[S.~1]{zhao2023survey}). Mit steigender Anzahl an Parametern, wachsender Menge an Trainingsdaten und höherem Rechenaufwand nimmt die Leistungsfähigkeit der \acp{LLM} dabei systematisch zu (vgl. \cite[S.~3]{kaplan2020scaling}).

Um Text verarbeiten zu können, zerlegt ein \ac{LLM} diesen zunächst in kleinere Einheiten, sogenannte Tokens. Ein Token kann ein ganzes Wort, einen Wortbestandteil oder ein einzelnes Zeichen umfassen (vgl. \cite[S.~1715~f.]{sennrich2016neural}). Jedes Token wird anschließend in einen realwertigen Vektor in einem hochdimensionalen Raum überführt, ein sogenanntes Embedding. Da ähnliche Wörter ähnliche Vektoren erhalten, kann das Modell semantische Beziehungen zwischen Wörtern abbilden (vgl. \cite[S. 1137 ff.]{bengio2003neural}). Wie viele Tokens ein Modell gleichzeitig verarbeiten kann, wird durch das sogenannte Kontextfenster begrenzt (vgl. \cite[S.~157~f.]{liu2024lost}). Die Ausgabe erzeugt ein \ac{LLM} autoregressiv. Das bedeutet, dass der Text Token für Token entsteht. Jedes neue Token wird auf Basis der Eingabe und aller bereits erzeugten Tokens vorhergesagt, an die bisherige Folge angehängt und fließt so in die nächste Vorhersage ein (vgl. \cite[S.~1138]{bengio2003neural}).

Auf dieser Grundlage können \acp{LLM} eine Vielzahl sprachbasierter Aufgaben bearbeiten, etwa Fragen beantworten, Texte übersetzen und zusammenfassen oder Informationen extrahieren (vgl. \cite[S.~1]{zhao2023survey}). Welche Aufgabe gelöst werden soll, erfährt das Modell über das sogenannte In-Context-Learning. Dabei wird die Aufgabe direkt in der Eingabe (Prompt) durch eine Anweisung in natürlicher Sprache beschrieben. Das Modell erzeugt die Lösung dann, indem es diese Eingabe Token für Token fortsetzt, ohne dass seine Parameter verändert werden. Dadurch lassen sich \acp{LLM} ohne aufgabenspezifisches Training flexibel für unterschiedliche Anwendungsfälle einsetzen (vgl. \cite[S.~1877~f.]{brown2020language}). Multimodale \acp{LLM} erweitern diese Fähigkeiten, indem sie Sprachmodelle mit Modellen zur Bildverarbeitung kombinieren. So können auch visuelle Eingaben interpretiert und beispielsweise in Bildern enthaltene Texte ohne vorgelagerte Texterkennung ausgewertet werden (vgl. \cite[S.~1]{yin2024survey}).

Diesen Fähigkeiten stehen jedoch wesentliche Einschränkungen gegenüber. So neigen \acp{LLM} zu Halluzinationen, also zu der Generierung von Inhalten, die plausibel erscheinen, inhaltlich jedoch falsch oder nicht belegbar sind (vgl. \cites[S.~2]{ji2023hallucination}[S.~1~f.]{huang2025survey}). Zudem arbeiten \acp{LLM} nicht vollständig deterministisch. Da sie ihre Ausgabe auf Basis von Wahrscheinlichkeiten erzeugen, sind identische Ergebnisse bei wiederholter Verarbeitung nicht garantiert (vgl. \cite[S.~5, 12]{huang2025survey}). Selbst identische Eingaben bei identischer Konfiguration können zu unterschiedlichen Ausgaben führen (vgl. \cite[S.~1]{ouyang2025empirical}). In Anwendungsfeldern mit hohen Anforderungen an Korrektheit und Nachvollziehbarkeit wie dem Internal Audit müssen die Ergebnisse eines \ac{LLM} daher fachlich überprüft und ihre Konsistenz gesondert untersucht werden.

\section{Prompt Engineering}
\label{sec:promptengineering}

Ein Prompt ist die Eingabe, mit der die Ausgabe eines \ac{LLM} gesteuert wird. Er kann unter anderem eine Anweisung zur eigentlichen Aufgabe, eine Rolle, zusätzliche Informationen, Beispiele sowie Vorgaben zum Ausgabeformat enthalten (vgl. \cite[S.~5~f.]{schulhoff2025promptreportsystematicsurvey}). Prompt Engineering bezeichnet den iterativen Prozess, einen Prompt auf Daten anzuwenden, die Ergebnisse zu bewerten und den Prompt auf dieser Grundlage anzupassen (vgl. \cite[S.~7]{schulhoff2025promptreportsystematicsurvey}).

Grundlegend wird zwischen Zero-Shot-Prompting, bei dem das Modell nur eine Beschreibung der Aufgabe erhält, und Few-Shot-Prompting mit zusätzlichen Lösungsbeispielen unterschieden (vgl. \cite[S.~1879]{brown2020language}). Beispiele können die Ergebnisse verbessern, belegen aber Platz im Kontextfenster (vgl. \cite[S.~10~f.]{schulhoff2025promptreportsystematicsurvey}). Ohne Beispiele ist das Modell auf eine klare Anweisung angewiesen. Diese kann etwa durch die Zuweisung einer Rolle (Role Prompting) und die ausdrückliche Zerlegung komplexer Aufgaben in einfachere Teilschritte verbessert werden (vgl. \cite[S.~11, 13]{schulhoff2025promptreportsystematicsurvey}). Bereits kleine Änderungen am Prompt, etwa andere Trennzeichen oder eine geänderte Reihenfolge, können die Ergebnisse deutlich beeinflussen (vgl. \cite[S.~31]{schulhoff2025promptreportsystematicsurvey}). Ein Prompt sollte daher am konkreten Anwendungsfall erprobt werden.

\section{KI-Agenten}

Unter einem autonomen Agenten versteht man ein System, das seine Umgebung wahrnimmt und in dieser handeln kann, um ein konkretes Ziel zu erreichen (vgl. \cite[S. 1]{wang2024survey}). Agenten basieren dabei auf \acp{LLM}, welche das Verhalten der Agenten steuern. Der Agent wird dabei mit weiteren Fähigkeiten ergänzt. Dazu gehören eine im Prompt festgelegte Rolle, ein Gedächtnis, eine Planungskomponente sowie die Fähigkeit, Aktionen auszuführen und externe Werkzeuge zu nutzen (vgl. \cite[S. 6]{wang2024survey}).

Wie selbstständig ein Agent arbeitet, kann sich dabei deutlich unterscheiden. Daher lassen sich einzelne \ac{KI}-Agenten von sogenannter Agentic AI abgrenzen. Einzelne Agenten erledigen klar abgegrenzte Aufgaben und greifen dafür auf Werkzeuge zurück. Bei Agentic AI arbeiten dagegen mehrere spezialisierte Agenten zusammen, teilen Aufgaben selbstständig untereinander auf und nutzen ein einheitliches Gedächtnis (vgl. \cite[S. 5 f.]{sapkota2026agenticai}).

Gesteuert wird das Verhalten eines solchen Agenten vor allem über Anweisungen, die dauerhaft im Prompt hinterlegt sind. Sie legen Rolle, Aufgabe und Verhalten des Agenten fest (vgl. \cite[S. 6]{wang2024survey}). Techniken des Prompt Engineerings lassen sich daher direkt auf die Gestaltung eines Agenten übertragen. Die in dieser Arbeit erstellten Agenten zählen zu \ac{KI}-Agenten und gelten nicht als Agentic AI, da sie eine klar abgegrenzte Aufgabe ausführen (vgl. \cite[S. 1 ff.]{sapkota2026agenticai}).
